\documentclass{situationist-preprint}

\title{The Evidential Weight of Religious Experience}
\msid{REL-21-a}
\msfamily{REL-21}
\msdiscipline{Philosophy of Religion}
\mstype{Principal manuscript}
\mskeywords{religious experience; perceiving God; phenomenal conservatism; principle of credulity; cross-checking; evidential weight; religious diversity}

\begin{document}
\maketitle

\begin{abstract}
Many people report experiences in which God seems to be present to them. Do such experiences give evidence that God exists? We argue that they do, prima facie, but that their weight is low, because religious experience lacks the cross-checking procedures that give sense perception its weight. The first claim follows from phenomenal conservatism: if it seems to a subject that $p$, she has prima facie justification for believing $p$, and theistic experiences are seemings with theistic content. The second rests on a distinction between the existence of prima facie justification and its weight. We argue that the high weight of perception is not conferred by its seemings alone but by procedures that test them: agreement among observers placed in the same conditions, convergence across senses, prediction, and a theory of the conditions under which perception errs. Religious experience has no comparable procedures, and Alston's explanation of their absence, that God is not at our disposal, explains it without supplying what is missing. We show further that reports of religious experience cannot be aggregated as independent testimonies, since their contents are shaped by shared traditions, and that the diversity of their contents leaves support mainly for a generic and weak conclusion. Religious experience is evidence, but not strong evidence.
\end{abstract}

\section{The problem}

The argument from religious experience holds that experiences in which God seems to be present are evidence that God exists, in the way that perceptual experiences of a tree are evidence that there is a tree. \citet{James1902} gave it its modern starting point, arguing that mystical states are authoritative for those who have them but impose no duty on others to accept them. Two systematic defences followed. \citet{Swinburne1979} proposed a principle of credulity, on which, in the absence of special considerations, if it seems to a subject that $x$ is present, then probably $x$ is present, together with a principle of testimony on which others' reports of their experiences are probably accurate. \citet{Alston1991} argued that Christian mystical practice is a socially established doxastic practice, like sense perception, and that engaging in it is rational in the absence of sufficient reasons for regarding it as unreliable.

Critics have pressed three points. \citet{Martin1959} and \citet{Gale1991} argued that religious experience lacks the tests by which perceptual claims are checked, so that it is not cognitive at all, or at least not perceptual. \citet{Rowe1982} argued that the principle of credulity is undermined in the religious case by the absence of conditions under which an experience of God's absence would be expected. A third objection appeals to diversity: experiences in different traditions have incompatible contents, and \citet{Hick1989} concluded that they are responses to a single ultimate reality filtered through culturally shaped concepts, a conclusion that \citet{Katz1978} reached by a different route, arguing that there is no unmediated mystical experience.

We take an intermediate position. Against the critics, religious experience is evidence, since seemings confer prima facie justification. Against Swinburne and Alston, its weight is low, since what gives perception its weight is not the seemings alone but the procedures that check them, and religious experience lacks comparable procedures.

\section{Framework}

\begin{principle}[Phenomenal conservatism]\label{p:pc}
If it seems to a subject that $p$, she has prima facie justification for believing that $p$.
\end{principle}

\ref{p:pc} is the conclusion of EPI-19, and is defended by \citet{Huemer2007}. Swinburne's principle of credulity is its probabilistic counterpart.

\begin{definition}[Religious experience]\label{d:rel}
A \emph{theistic experience} is an experience in which it seems to the subject that God is present, or that God is acting toward her, where the seeming is not merely a belief or an inference from other experiences.
\end{definition}

The definition covers what \citet{Alston1991} calls the direct perception of God and excludes interpretations of ordinary experiences, such as the sense that nature is designed. It does not require sensory content: Alston's reports are typically of a non-sensory awareness of a presence.

\begin{definition}[Weight]\label{d:weight}
The \emph{weight} of a class of experiences, as evidence for its contents, is the degree to which a seeming of the class raises the probability of its content, on the background information available to an agent who assesses it.
\end{definition}

Prima facie justification and weight are distinct. A seeming may give some justification, which stands unless defeated, while raising the probability of its content only slightly. Phenomenal conservatism fixes the existence of prima facie justification; it does not fix its weight.

\begin{definition}[Cross-checking procedure]\label{d:check}
A \emph{cross-checking procedure} for a class of experiences is a procedure, other than having a further experience of the same kind by the same subject, whose outcome can disagree with a given experience of the class and whose disagreement counts against it. The main cross-checks of sense perception are: (i) agreement of other observers placed in the same conditions; (ii) convergence of different sense modalities; (iii) successful prediction of further experiences; and (iv) a theory of the conditions under which the faculty errs.
\end{definition}

\begin{principle}[Calibration]\label{p:calib}
The weight of a class of seemings beyond the minimal weight conferred by \ref{p:pc} is proportional to the extent to which its reliability has been confirmed by cross-checking procedures. In the absence of such procedures, its weight is low.
\end{principle}

\ref{p:calib} is the principle we use most. It says that seemings earn weight by being tested, and it applies to every class of seemings: to perception, to memory, to the intuitions of experts in a field, and to clinical judgement. A radiologist's impression that an image shows a tumour weighs much more than a novice's, though both are seemings, because hers has been checked against biopsies.

\section{The argument}

\begin{proposition}[Prima facie evidence]\label{p:prima}
Theistic experiences give their subjects prima facie justification for believing that God is present, and hence that God exists.
\end{proposition}

\begin{proof}[Argument]
A theistic experience is a seeming that God is present (Definition~\ref{d:rel}). By \ref{p:pc}, its subject has prima facie justification for believing its content. That God is present entails that God exists.
\end{proof}

Proposition~\ref{p:prima} answers the critics who deny that religious experience is evidence at all. Their arguments establish at most that the experience is defeasible and lacks certain tests, not that it confers no justification. A principle on which seemings justify only when they are of a privileged kind would need to say which kind, and the usual answer, sensory seemings, is too narrow: memory and rational intuition also justify without being sensory.

\begin{lemma}[Perception's weight is earned]\label{l:earned}
The high weight of sense perception depends on its cross-checking procedures, and not on \ref{p:pc} alone.
\end{lemma}

\begin{proof}[Argument]
Consider a perceptual seeming in conditions where all four checks of Definition~\ref{d:check} are unavailable: a faint light seen once, by one observer, in conditions of which nothing is known, which no other sense can confirm and from which nothing follows. Such a seeming confers prima facie justification, but we give it little weight, and we are right to. Its weight rises as the checks become available: when others see it too, when it is photographed, when its position predicts where it will be seen next. The weight of perception in general is therefore the product of its seemings and of a long history of successful cross-checking, which has also given us a theory of illusions and of the conditions under which they occur. By \ref{p:calib}, this is as it should be.
\end{proof}

\begin{lemma}[Absence of cross-checks]\label{l:absent}
Theistic experience has no cross-checking procedures comparable to (i)--(iv).
\end{lemma}

\begin{proof}[Argument]
(i) There are no conditions such that anyone placed in them, with normal capacities, would have a theistic experience if God were present; most people who pray, fast or meditate in the prescribed ways do not have one. (ii) There is no second modality through which the presence of God could be independently confirmed. (iii) Theistic experiences yield few predictions about further experience that could fail: the predictions offered by mystical practice, such as growth in patience and humility, are equally expected on naturalistic accounts of the practice. (iv) There is no theory of the conditions under which the faculty errs that is independent of the practice's own doctrine.
\end{proof}

\citet{Alston1991} accepts much of Lemma~\ref{l:absent} and argues that the demand for such checks is a form of epistemic imperialism: it imposes the standards of one doxastic practice on another whose object, a free personal being who is not at our disposal, would not be expected to satisfy them. We grant the explanation. If God exists and chooses when to be manifest, condition (i) cannot be expected to hold. But an explanation of why the checks are absent does not supply what they would have supplied. By \ref{p:calib}, weight depends on confirmed reliability, and a reason why reliability cannot be confirmed is not a confirmation of it. The point is general: an observer whose instrument cannot be calibrated for principled reasons has a good excuse, and a weak measurement.

\begin{theorem}\label{t:main}
Theistic experience is prima facie evidence that God exists, and its weight is low.
\end{theorem}

\begin{proof}[Argument]
The first conjunct is Proposition~\ref{p:prima}. By Lemma~\ref{l:absent}, theistic experience lacks cross-checking procedures comparable to those of perception, which by Lemma~\ref{l:earned} are what give perception its weight. By \ref{p:calib}, its weight is low.
\end{proof}

\section{Aggregation and diversity}

Two further considerations bear on how much the evidence of many experiences adds up to.

\begin{proposition}[Limited aggregation]\label{p:agg}
Reports of theistic experience by many subjects do not combine as independent testimonies.
\end{proposition}

\begin{proof}[Argument]
REL-07 shows that Hume's maxim can be met by a sufficient number of testimonies that are independent conditional on the truth of the report. The condition matters: independent witnesses multiply likelihood ratios, correlated ones do not. Reports of theistic experience are not independent conditional on either hypothesis. Their contents, and the occasions on which they occur, are shaped by shared practices, texts and expectations \citep{Katz1978}, so that a further report from within a tradition is to a large extent predicted by the earlier ones whether or not God exists. The number of reports is therefore much less informative than the number of independent witnesses would be.
\end{proof}

\begin{proposition}[Diversity]\label{p:div}
The diversity of the contents of religious experience leaves stronger support for generic contents than for the specific contents of any one tradition.
\end{proposition}

\begin{proof}[Argument]
Experiences across traditions present a personal God, an impersonal absolute, or an emptiness that is neither. If each is prima facie evidence for its content, the specific contents defeat one another in part, as conflicting perceptual reports do. What they share, the seeming presence of a reality that transcends the ordinary, receives the support of all of them, but this generic content is much weaker than theism and is consistent with several non-theistic hypotheses \citep{Hick1989}.
\end{proof}

Propositions~\ref{p:agg} and~\ref{p:div} reinforce Theorem~\ref{t:main} from a different direction. The weight of a single experience is low because it cannot be checked, and the weight of many experiences together is not proportionally greater, because they are not independent and they disagree in what they present.

\section{Objections}

\begin{objection}[Epistemic circularity]
Sense perception's cross-checks are themselves perceptual: we check one observation with another. Its reliability cannot be established without circularity, and Alston showed that the same holds of every basic doxastic practice. To demand of religious experience a non-circular vindication is a double standard.
\end{objection}

\begin{reply}
We demand no non-circular vindication. Cross-checks internal to a practice are worth having if they could fail, and perceptual cross-checks often do: observers disagree, predictions fail, and illusions are detected. A practice whose internal checks can fail and mostly do not has earned weight from within, even if its reliability cannot be shown from without. Lemma~\ref{l:absent} says that theistic experience lacks internal checks of this kind, not that it lacks external vindication.
\end{reply}

\begin{objection}[Credulity]
By Swinburne's principle of credulity, a seeming that God is present makes it probable that God is present unless there are special considerations against it. No such considerations have been identified, so the weight is high \citep{Swinburne1979}.
\end{objection}

\begin{reply}
The principle of credulity, like \ref{p:pc}, establishes a default. Its \enquote{probably} cannot mean the same thing for every class of seemings, since seemings differ in weight, as the radiologist and the novice show. Swinburne's own special considerations include the subject's being in conditions previously shown to give unreliable experiences, which presupposes a theory of when the faculty errs. Where there is no such theory, the default is not raised to high probability; it remains a default, of low weight.
\end{reply}

\begin{objection}[First-person weight]
The mystic's own experience may be overwhelming: vivid, stable, and transformative. Whatever its weight for others, for her it is decisive.
\end{objection}

\begin{reply}
Vividness and stability increase the strength of a seeming, and so of the prima facie justification it confers. They are not calibration: vivid and stable experiences are also produced by fever and by drugs. The subject's position differs from an observer's in that she has the seeming itself rather than a report of it, and REL-11 shows that in a peer disagreement steadfastness can be rational for at most one party; to the extent that the experience is reported and weighed, its residue for the subject alone is limited. That residue may rationally move her, and we do not deny it. It does not make the evidence weighty by the standard of \ref{p:calib}.
\end{reply}

\begin{objection}[Debunking]
Cognitive science explains religious experience by mechanisms, such as hypersensitive agency detection, that would produce it whether or not God exists. Its weight is therefore zero.
\end{objection}

\begin{reply}
An etiology that would produce the experience whether or not its content is true undercuts it only if it is known to be the complete etiology. A theist may hold that God acts through natural mechanisms, as the perceptual system works through retinal processes. The debunking explanations are partial and contested, and their effect is to lower the weight of the evidence further, not to cancel it. This is consistent with Theorem~\ref{t:main}.
\end{reply}

\section{What the result does not establish}

The theorem does not imply that God does not exist, or that theistic belief is irrational. A belief can be rational on several pieces of evidence each of low weight, and the arguments for and against theism examined in REL-01, REL-02 and REL-10 are assessed on their own terms. Nor does it imply that religious commitment requires stronger evidence than religious experience provides: REL-17 argues that faith, as a practical commitment, is rational where the evidence underdetermines the question and the stakes are practical, and a body of evidence that includes religious experience of low weight is one under which the question may remain underdetermined. Finally, the result is not that religious experience could never acquire weight. If a procedure were found by which its reports could be tested and could fail, its weight would rise in proportion to their success.

\section{Conclusion}

Religious experience is evidence, since seemings justify prima facie, and theistic experiences are seemings. Its weight is low, because the weight of a class of seemings is earned by procedures that check it, and theistic experience has no procedures comparable to those of perception. The explanation of their absence does not substitute for them. Reports from many subjects do not combine as independent testimonies, and the diversity of their contents leaves support chiefly for a generic conclusion. The argument from religious experience, taken alone, gives some reason to believe that God exists, but not much.

\begin{thebibliography}{99}\small
\bibitem[Alston(1991)]{Alston1991} Alston, W.~P. (1991). \emph{Perceiving God: The Epistemology of Religious Experience}. Cornell University Press.
\bibitem[Gale(1991)]{Gale1991} Gale, R.~M. (1991). \emph{On the Nature and Existence of God}. Cambridge University Press.
\bibitem[Hick(1989)]{Hick1989} Hick, J. (1989). \emph{An Interpretation of Religion}. Macmillan.
\bibitem[Huemer(2007)]{Huemer2007} Huemer, M. (2007). Compassionate phenomenal conservatism. \emph{Philosophy and Phenomenological Research}, 74(1), 30--55.
\bibitem[James(1902)]{James1902} James, W. (1902). \emph{The Varieties of Religious Experience}. Longmans, Green.
\bibitem[Katz(1978)]{Katz1978} Katz, S.~T. (1978). Language, epistemology, and mysticism. In S.~T. Katz (ed.), \emph{Mysticism and Philosophical Analysis}. Oxford University Press.
\bibitem[Martin(1959)]{Martin1959} Martin, C.~B. (1959). \emph{Religious Belief}. Cornell University Press.
\bibitem[Rowe(1982)]{Rowe1982} Rowe, W.~L. (1982). Religious experience and the principle of credulity. \emph{International Journal for Philosophy of Religion}, 13(2), 85--92.
\bibitem[Swinburne(1979)]{Swinburne1979} Swinburne, R. (1979). \emph{The Existence of God}. Clarendon Press.
\end{thebibliography}

\end{document}
